% Options for packages loaded elsewhere
\PassOptionsToPackage{unicode}{hyperref}
\PassOptionsToPackage{hyphens}{url}
\PassOptionsToPackage{space}{xeCJK}
%
\documentclass[
  aps,
  onecolumn,
  11pt]{revtex4-2}
\usepackage{amsmath,amssymb}
\usepackage{iftex}
\ifPDFTeX
  \usepackage[T1]{fontenc}
  \usepackage[utf8]{inputenc}
  \usepackage{textcomp} % provide euro and other symbols
\else % if luatex or xetex
  \usepackage{unicode-math} % this also loads fontspec
  \defaultfontfeatures{Scale=MatchLowercase}
  \defaultfontfeatures[\rmfamily]{Ligatures=TeX,Scale=1}
\fi
\usepackage{lmodern}
\ifPDFTeX\else
  % xetex/luatex font selection
  \setmainfont[]{Times New Roman}
  \setmonofont[]{Menlo}
  \ifXeTeX
    \usepackage{xeCJK}
    \setCJKmainfont[]{Songti SC}
          \fi
  \ifLuaTeX
    \usepackage[]{luatexja-fontspec}
    \setmainjfont[]{Songti SC}
  \fi
\fi
% Use upquote if available, for straight quotes in verbatim environments
\IfFileExists{upquote.sty}{\usepackage{upquote}}{}
\IfFileExists{microtype.sty}{% use microtype if available
  \usepackage[]{microtype}
  \UseMicrotypeSet[protrusion]{basicmath} % disable protrusion for tt fonts
}{}
\makeatletter
\@ifundefined{KOMAClassName}{% if non-KOMA class
  \IfFileExists{parskip.sty}{%
    \usepackage{parskip}
  }{% else
    \setlength{\parindent}{0pt}
    \setlength{\parskip}{6pt plus 2pt minus 1pt}}
}{% if KOMA class
  \KOMAoptions{parskip=half}}
\makeatother
\usepackage{xcolor}
\usepackage{color}
\usepackage{fancyvrb}
\newcommand{\VerbBar}{|}
\newcommand{\VERB}{\Verb[commandchars=\\\{\}]}
\DefineVerbatimEnvironment{Highlighting}{Verbatim}{commandchars=\\\{\}}
% Add ',fontsize=\small' for more characters per line
\newenvironment{Shaded}{}{}
\newcommand{\AlertTok}[1]{\textcolor[rgb]{1.00,0.00,0.00}{\textbf{#1}}}
\newcommand{\AnnotationTok}[1]{\textcolor[rgb]{0.38,0.63,0.69}{\textbf{\textit{#1}}}}
\newcommand{\AttributeTok}[1]{\textcolor[rgb]{0.49,0.56,0.16}{#1}}
\newcommand{\BaseNTok}[1]{\textcolor[rgb]{0.25,0.63,0.44}{#1}}
\newcommand{\BuiltInTok}[1]{\textcolor[rgb]{0.00,0.50,0.00}{#1}}
\newcommand{\CharTok}[1]{\textcolor[rgb]{0.25,0.44,0.63}{#1}}
\newcommand{\CommentTok}[1]{\textcolor[rgb]{0.38,0.63,0.69}{\textit{#1}}}
\newcommand{\CommentVarTok}[1]{\textcolor[rgb]{0.38,0.63,0.69}{\textbf{\textit{#1}}}}
\newcommand{\ConstantTok}[1]{\textcolor[rgb]{0.53,0.00,0.00}{#1}}
\newcommand{\ControlFlowTok}[1]{\textcolor[rgb]{0.00,0.44,0.13}{\textbf{#1}}}
\newcommand{\DataTypeTok}[1]{\textcolor[rgb]{0.56,0.13,0.00}{#1}}
\newcommand{\DecValTok}[1]{\textcolor[rgb]{0.25,0.63,0.44}{#1}}
\newcommand{\DocumentationTok}[1]{\textcolor[rgb]{0.73,0.13,0.13}{\textit{#1}}}
\newcommand{\ErrorTok}[1]{\textcolor[rgb]{1.00,0.00,0.00}{\textbf{#1}}}
\newcommand{\ExtensionTok}[1]{#1}
\newcommand{\FloatTok}[1]{\textcolor[rgb]{0.25,0.63,0.44}{#1}}
\newcommand{\FunctionTok}[1]{\textcolor[rgb]{0.02,0.16,0.49}{#1}}
\newcommand{\ImportTok}[1]{\textcolor[rgb]{0.00,0.50,0.00}{\textbf{#1}}}
\newcommand{\InformationTok}[1]{\textcolor[rgb]{0.38,0.63,0.69}{\textbf{\textit{#1}}}}
\newcommand{\KeywordTok}[1]{\textcolor[rgb]{0.00,0.44,0.13}{\textbf{#1}}}
\newcommand{\NormalTok}[1]{#1}
\newcommand{\OperatorTok}[1]{\textcolor[rgb]{0.40,0.40,0.40}{#1}}
\newcommand{\OtherTok}[1]{\textcolor[rgb]{0.00,0.44,0.13}{#1}}
\newcommand{\PreprocessorTok}[1]{\textcolor[rgb]{0.74,0.48,0.00}{#1}}
\newcommand{\RegionMarkerTok}[1]{#1}
\newcommand{\SpecialCharTok}[1]{\textcolor[rgb]{0.25,0.44,0.63}{#1}}
\newcommand{\SpecialStringTok}[1]{\textcolor[rgb]{0.73,0.40,0.53}{#1}}
\newcommand{\StringTok}[1]{\textcolor[rgb]{0.25,0.44,0.63}{#1}}
\newcommand{\VariableTok}[1]{\textcolor[rgb]{0.10,0.09,0.49}{#1}}
\newcommand{\VerbatimStringTok}[1]{\textcolor[rgb]{0.25,0.44,0.63}{#1}}
\newcommand{\WarningTok}[1]{\textcolor[rgb]{0.38,0.63,0.69}{\textbf{\textit{#1}}}}
\setlength{\emergencystretch}{3em} % prevent overfull lines
\providecommand{\tightlist}{%
  \setlength{\itemsep}{0pt}\setlength{\parskip}{0pt}}
\setcounter{secnumdepth}{-\maxdimen} % remove section numbering
\ifLuaTeX
  \usepackage{selnolig}  % disable illegal ligatures
\fi
\IfFileExists{bookmark.sty}{\usepackage{bookmark}}{\usepackage{hyperref}}
\IfFileExists{xurl.sty}{\usepackage{xurl}}{} % add URL line breaks if available
\urlstyle{same}
\hypersetup{
  pdftitle={Mathieu Stability User Guide},
  pdfauthor={TrapLab documentation},
  hidelinks,
  pdfcreator={LaTeX via pandoc}}

\title{Mathieu Stability User Guide}
\author{TrapLab documentation}
\date{2026-10-24}

\begin{document}
\maketitle

This guide explains the separate Mathieu stability explorer on the
website and in the installed companion application. It calculates ideal
bounded motion in one or more Mathieu bands. A stable point is not a
measured transmission, a cooling efficiency or an electrode-clearance
result. The classic Keep the ions game remains a separate trajectory
demonstration.

\subsection{First calculation}\label{first-calculation}

\begin{enumerate}
\def\labelenumi{\arabic{enumi}.}
\tightlist
\item
  Open Games, go to Keep the ions, and select Mathieu bands. The same
  independent explorer is available on Lab. Choose \textbf{Linear RFQ},
  \textbf{Direct point}, three bands, horizontal limit 2 and vertical
  limits −1 to 1.
\item
  Enter \(a=0\) and \(q=0.4\). Both transverse planes should report
  stable in the first band. Change \(a\) or \(q\) and read each plane's
  status rather than judging the shaded picture alone.
\item
  Switch to voltage input. Use ion mass 100 u, charge 1, \(U=0\) V,
  \(V=122\) V, \(r_0=10\) mm and \(f=1.3\) MHz. The calculation gives
  \(a=0\) and \(q\simeq0.07057237\). The shipped voltage inputs have
  \(U=3\) V instead; those give \(a\simeq0.00347077\). Do not assume
  that every default operating point is stable in both planes.
\item
  Change the three section settings independently. Cooler defaults are 3
  V, 122 V, 10 mm and 1.3 MHz; mass-filter and buncher defaults are each
  3 V, 244 V, 3.75 mm and 3 MHz. Read their separate lines and ion
  markers. These overlays belong to the linear RFQ geometry.
\item
  Open \textbf{Figure style} to change text, reference marker diameter,
  line width, raster width or filename prefix. These controls change
  drawing only. Save settings JSON, change an input, reopen the saved
  JSON, and check that the inputs, axes and style return.
\end{enumerate}

The Species field and \textbf{Fill ion mass} explicitly obtain an ion
mass from the existing isotope catalogue. Manual mass and charge remain
editable; merely typing a species does not perform an automatic fill.
Record the returned source metadata when using an evaluated mass.

\subsection{Equation and voltage
convention}\label{equation-and-voltage-convention}

The independent variable and standard equation are

\[\tau=\frac{\Omega t}{2},\qquad \Omega=2\pi f,\]

\[\frac{d^2w}{d\tau^2}+\left[a-2q\cos(2\tau)\right]w=0.\]

For this linear-rod convention,

\[a=\frac{8zeU}{m r_0^2\Omega^2},\qquad q=\frac{4zeV}{m r_0^2\Omega^2}.\]

Here \(z\) is the positive integer charge state, \(e\) is the elementary
charge, and \(m\) is the actual ion mass in kg. The UI accepts mass in
u, radius in mm and frequency in MHz and converts them to SI. \(U\) is
signed DC voltage and \(V\) is the rod-to-ground RF peak. The
differential RF peak is \(2V\); RMS or peak-to-peak values must be
converted before entry. The second transverse plane uses

\[a_y=-a_x,\qquad q_y=-q_x.\]

Both planes must be stable for ideal transverse confinement. Reducing
the rod radius or frequency increases the magnitude of the parameters
quadratically. Charge must be an integer from 1 to 100; mass, radius and
frequency must be positive and finite, and RF peak must be nonnegative.

\subsection{Bands and the separate three dimensional
geometry}\label{bands-and-the-separate-three-dimensional-geometry}

For nonzero real \(q\), the displayed one-dimensional bands are indexed
by \(n=0,\ldots,N-1\) and satisfy

\[a_n(|q|)<a<b_{n+1}(|q|).\]

The UI supports one to six bands and \(|q|\leq12\). A boundary is
reported separately. Above the calculated bands or outside the covered
domain, the status is outside coverage, not a demonstrated instability.
At \(q=0\), positive \(a\) is treated analytically as stable;
\((a,q)=(0,0)\) is a boundary.

The \textbf{3D Paul} view accepts direct axial parameters only. It uses

\[a_r=-\frac{a_z}{2},\qquad q_r=-\frac{q_z}{2}.\]

Its radial boundary curves in axial coordinates therefore are

\[-2a_n(|q_z|/2),\qquad -2b_{n+1}(|q_z|/2).\]

The interface does not infer an electrode voltage convention for this
geometry. The three linear-section lines and points are hidden in this
view. Do not transfer the linear-rod voltage formula to an arbitrary
three-dimensional electrode layout.

\subsection{Numerical checks and saved
files}\label{numerical-checks-and-saved-files}

The browser computes characteristic values from parity-specific
symmetric Hill matrices using Sturm bisection. Native computation uses
installed SciPy Mathieu characteristic functions. An independent RK4
integration over one coefficient period constructs a monodromy matrix
\(M\). Its diagnostics include

\[\det M\simeq1,\qquad |\operatorname{tr}M|<2,\]

\[\Delta_{\mathrm{tr}}=|\operatorname{tr}M_{1024}-\operatorname{tr}M_{512}|.\]

The strict trace inequality describes the interior stable case. Near
boundaries, inspect the characteristic-value classification and
convergence diagnostic. Failure of the numerical diagnostic to converge
is unresolved. The shaded diagram samples 241 horizontal points, so
narrow higher-band regions can be visually undersampled; the selected
point uses characteristic values directly.

\textbf{SVG} saves vector curves and text. \textbf{PNG} uses the
selected raster width. \textbf{CSV} saves the sampled band data;
\textbf{Save settings JSON} preserves geometry, inputs, section
parameters, axes and style for reopening. Browser files go to the
browser's configured download location. The existing shared image-PDF
facility, where exposed, is a raster figure PDF; it is distinct from the
vector SVG and from compiling this guide's TeX source.

\subsection{Installed companion
workflow}\label{installed-companion-workflow}

In the separately installed companion folder, open
\texttt{Individual\_tools/13\_Mathieu\_stability.command}, or run

\begin{Shaded}
\begin{Highlighting}[]
\ExtensionTok{./Run\_private.command}\NormalTok{ traplab.py gui }\AttributeTok{{-}{-}standalone} \AttributeTok{{-}{-}page}\NormalTok{ mathieu\_stability}
\end{Highlighting}
\end{Shaded}

The integrated page is Mathieu stability multiple bands under Paul trap
and RFQ. It starts with a selected 133Cs ion and provides one selected
operating line, direct or voltage inputs, local axis font and marker
diameter, normal page-state persistence, and the existing Plot style
dock and Export plots actions. The existing \texttt{rfq\_sections} page
retains the native three-section workflow; the new native page does not
add three independent section lines. Native launch modules are portable
entry-point labels, not website downloads.

Web SI u is \(1.66053906660\times10^{-27}\) kg from its retained
CODATA2018 profile; the native profile uses
\(1.66053906892\times10^{-27}\) kg from CODATA2022. Small parameter
differences are expected. No uncertainty propagation, fringe fields,
collisions, space charge, electrode losses or hardware acceptance is
included.

\subsection{Recovery and evidence}\label{recovery-and-evidence}

For a blank or invalid input, restore a finite value in the stated
range. If a requested point is outside coverage, reduce the domain or
increase the band count within the supported limits; do not relabel it
stable. If a species fill fails, retain its exact isotope/state or
explicitly enter an ion mass. Reopen settings of this explorer rather
than a legacy diagram file.

Source checks covered 150 characteristic values at six selected \(q\)
values and orders through 12 against SciPy 1.10.1, with maximum absolute
difference \(2.85\times10^{-14}\). Nine focused web tests and five
native offscreen tests passed. These bounded checks do not establish
browser download acceptance, native Cocoa export acceptance or
experimental performance.

\subsection{Primary references}\label{primary-references}

\begin{itemize}
\tightlist
\item
  NIST Digital Library of Mathematical Functions, version 1.2.8,
  \href{https://dlmf.nist.gov/28.2}{28.2 Definitions and Basic
  Properties} and \href{https://dlmf.nist.gov/28.17}{28.17 Stability},
  checked 24 October 2026.
\item
  SciPy official documentation,
  \href{https://docs.scipy.org/doc/scipy/reference/generated/scipy.special.mathieu_a.html}{mathieu\_a}
  and
  \href{https://docs.scipy.org/doc/scipy/reference/generated/scipy.special.mathieu_b.html}{mathieu\_b}.
  The installed native validation version is 1.10.1; current online
  documentation may describe a later release.
\item
  Evaluated mass provenance, when filling an ion, is retained from
  \href{https://www-nds.iaea.org/amdc/ame2020/mass_1.mas20.txt}{AME2020}
  and
  \href{https://www-nds.iaea.org/amdc/ame2020/nubase_4.mas20.txt}{NUBASE2020}.
  Numerical conventions reuse the existing companion
  \texttt{paul\_mathieu} and \texttt{rfq\_sections} implementations.
\end{itemize}

\end{document}
